% ---- pg-586 （书页 574） ----

\begin{mvfigure}
\mvinc{pg-586}{0.287,0.131,0.714,0.413}
\phantomsection\label{fig:12.29}\mvcap{图 12.29\quad 森林数据的双重信息图}
\end{mvfigure}

这两个点图给出非常相似的信息. 场所 1 至 5 倾向于与橡树类相联系, 而场所 7 至 10 倾向于与菩提树、铁树和糖枫相联系. 美国榆树分布在大部分场所. 由不同树种区分的场所几乎是连续的. \hfill$\blacksquare$

\mvsub{补充 12A\quad 数据挖掘}

\mvsub{引言}

在传统的统计方法中, 关于大约 50 个变量的 10 000 个观测意味着一组非常大的样本. 而今天, 基于计算机的数据存储设备可以存储万亿字节的数据. 对于某些组织, 处理的数据在最近几十年里增长了上十万倍. 电信、银行、制药以及集装箱航运等行业中都有一些公司有巨大的数据库. 考虑下面的例子. 假定在议会图书馆中的 1 700 万本图书平均相当于 MS Word 格式中大小为 1MB 的文本, 则将这些印刷材料输入计算机数据库将占用 17 万亿字节的磁盘空间. 联合邮递服务公司(United Parcel Service, UPS)有一个 17 万亿字节的数据库来管理其业务.

就我们的目标来说, 数据挖掘(data mining)意味着在极其巨大的数据内发现模式或者关系. 也就是说, 数据挖掘是要从数值信息的大山中挖掘出少量知识的金块. 从商业角度说, 这些知识的金块代表“活的”信息, 有助于发现获得竞争优势的方法.

数据挖掘需要用适合的软件和快速的计算机. 显而易见, 本书中介绍的各种数据分析技术与机器学习以及人工智能中开发的各种算法在数据挖掘中扮演着重要角色. 开发著名的统计软件包的公司现在提供包罗万象的数据挖掘程序.\footnote{SAS 公司的数据挖掘程序现在称为 Enterprise Miner. SPSS 的数据挖掘程序是 Clementine.} 此外, 一些特殊用途的程序, 如 CART, 在数据挖掘的应用中已经取得成功.

数据挖掘帮助确定用于处方药物的化学成分, 发现欺诈索赔和购买, 创建并维护个人客
